\subsection{赋范李代数}

定义 1。\(\mathbf{K}\) 上的赋范李代数是带有范数的李代数，满足

\begin{enumerate}
\def\labelenumi{(\arabic{enumi})}
\setcounter{enumi}{7}
\tightlist
\item
\end{enumerate}

\[
\| x + y \| \leqslant \sup (\| x \|, \| y \|)\tag{8}
\]

\[
\| [ x, y ] \| \leqslant \| x \|. \| y \|\tag{9}
\]

对 g 中所有 x、y。

在本段其余部分，\(\mathfrak{g}\) 表示完备赋范李代数。

对每个有限集 I，仿照 § 7 第 1 号定义从 \(u \mapsto \tilde{u}\) 到 \(\hat{\mathbf{L}}(\mathbf{I})\) 的连续李代数同态 \(\hat{\mathrm{P}}(\mathfrak{g}^{\mathrm{I}}; \mathfrak{g})\)。如 § 7 中所见，若 \(u = \sum_{\nu} u_{\nu}\)，其中对 \(u_{\nu} \in \mathbf{L}^{\nu}(\mathbf{I})\) 有 \(\nu \in \mathbf{N}^{\mathbf{I}}\)，则 \(\tilde{u} = \sum_{\nu} \tilde{u}_{\nu}\)，其中 \(\tilde{u}_{\nu}\) 是 § 2 第 4 号中定义的多项式映射 \((t_i)_{i \in \mathbf{I}} \mapsto u_{\nu}(t_i)\)。§ 7 第 1 号的复合公式 (2) 仍然成立。

